% Comandos personalizados

% \fechaPresentacion :: para escribir la fecha de presentación del trabajo
\newcommand{\fechaPresentacion}{\today}
% \unlp :: para escribir "Universidad Nacional de La Plata"
\newcommand{\unlp}{Universidad Nacional de La Plata\xspace}
% \facultad :: para escribir "Facultad de Informática"
\newcommand{\facultad}{Facultad de Informática}
% \cespi :: para escribir "CeSPI"
\newcommand{\cespi}{CeSPI}
% \direccionDesarrollo :: Para escribir "Dirección de Desarrollo del CeSPI"
\newcommand{\direccionDesarrollo}{Dirección de Desarrollo del \cespi\xspace}
% \tituloTrabajo :: Para escribir el título de la tesina
\newcommand{\tituloTrabajo}{Propuesta de rediseño de la nube de servicios de la UNLP}
% \tituloTrabajoDosLineas :: Para escribir el título de la tesina en dos líneas (carátula)
\newcommand{\tituloTrabajoDosLineas}{Propuesta de rediseño de la nube\\* de servicios de la UNLP}
% \miguelcarbone :: para escribir "Miguel Carbone"
\newcommand{\miguelcarbone}{Miguel Carbone}
% \carbonemiguel :: para escribir "Carbone, Miguel"
\newcommand{\carbonemiguel}{Carbone, Miguel}
% \nahuelcuesta :: para escribir "José Nahuel Cuesta Luengo"
\newcommand{\nahuelcuesta}{José Nahuel Cuesta Luengo}
% \cuestanahuel :: para escribir "Cuesta Luengo, José Nahuel"
\newcommand{\cuestanahuel}{Cuesta Luengo, José Nahuel}
% \cloud :: para escribir el nombre clave de la nueva nube
\newcommand{\cloud}{\textit{Cloud}}

% \eng{English expression} :: para denotar que "English expression" está en inglés
\newcommand{\eng}[1]{\textit{#1}}

% \caratula :: para generar la carátula de la tesina
\newcommand{\caratula}{
  \begin{center}
    \includegraphics{src/images/caratula/unlp.png}\\
    \huge{\unlp}\\
    \vspace{5mm}
    \huge{\facultad}\\
    \vspace{5mm}
    \large{Tesina de la Licenciatura}\\
    \vspace{15mm}
    \huge{\tituloTrabajoDosLineas}\\
    \vspace{10mm}
    \large{\textbf{\carbonemiguel} \\ \textbf{\cuestanahuel}}\\
    \vspace{20mm}
    \large{Directoras: Banchoff Tzancoff, Claudia y Queiruga, Claudia}\\
    \large{Asesor Profesional: Rodriguez, Christian Adrián}\\
    \vspace{20mm}
    \normalsize{\fechaPresentacion}\\
  \end{center}
}

% \checkmark :: para imprimir un check (tilde)
\def\checkmark{\tikz\fill[scale=0.4](0,.35) -- (.25,0) -- (1,.7) -- (.25,.15) -- cycle;}
